% ==================================================================
% id:       der-006
% topic:    נגזרות
% subtopic: הגדרת הנגזרת וגזירות
% solution: full
% ==================================================================
\textbf{מצאו את הנגזרות החד-צדדיות של הפונקציות הבאות בנקודה הנתונה ובדקו, האם הפונקציות גזירות בנקודה זו:}
\bE\bMs{2}
\item $x_0 = 2, \,\, f(x) = |x^2 - 5x + 6|$
\hint{פרקו את הביטוי שבתוך הערך המוחלט לגורמים, וכתבו את $f$ ללא ערך מוחלט משני צדי הנקודה.}
\item $x_0 = 1, \quad \ds{f(x) =\case{cc}{\frac{x-1}{4}(x+1)^2 & x \leq 1 \\|x| - 1 & x > 1}} $
\hint{חשבו כל נגזרת חד-צדדית לפי הנוסחה של הצד המתאים, ואל תשכחו לבדוק רציפות בנקודה.}
\eMs\eE

\ifsol
\begin{solbox}
\bE
\item
נציג
$f(x)=(x-2)(x-3)$
ולכן
נשים לב כי
$$.f(x) =\case{cc}{x^2 - 5x + 6 & x<2,\ x>3 \\-(x^2 - 5x + 6) & 2\leq x\leq3}$$
נסמן
$$.g(x)=x^2 - 5x + 6,\ h(x)=-(x^2 - 5x + 6)$$
מכאן
$$,f'_{-}(2)=g'(2)=-1,\ f'_{+}(2)=h'(2)=1$$
קיבלנו כי
$f'_{-}(2)\neq f'_{+}(2)$,
ולכן
$f$
לא גזירה ב-
$x=2$.

\item
נסמן
$$g(x)=\frac{x-1}{4}(x+1)^2,\ h(x)=|x|-1$$
נשים לב כי לכל
$x\geq1$
$h(x)=x-1$.
לכן
$$,f'_{-}(1)=g'_{-}(1)=1,\ f'_{+}(1)=h'_{+}(1)=1$$
קיבלנו כי
$f'_{-}(1)=f'_{+}(1)$
ובהוסף
$f$
רציפה
(בדקו זאת).
נסיק כי
$f$
גזירה ב-
$x=1$.

\eE
\end{solbox}
\fi
